\documentclass{nuspaper}

% Optional front-page overrides.
% The front page uses NUSGold for the title by default.
% Use \documentclass[yellowtitle]{nuspaper} for the brighter NUSYellow.
% To make it closer to the reference paper's crimson title, use for example:
% \definecolor{ReferenceCrimson}{HTML}{E01843}
% \colorlet{NUSFrontTitle}{ReferenceCrimson}

\title{NUSPaper: A Clean House Style for Research Preprints}
\shorttitle{NUSPaper: A Clean House Style for Research Preprints}
\date{2026-06-15}
\venue{Technical Report}
\version{Template preview v0.1}
\arxivstamp{arXiv:XXXX.XXXXXv1 [cs.LG] 15 Jun 2026}

\author[1,*]{First Author}
\author[1,*]{Second Author}
\author[2]{Third Author}
\author[1,\dag]{Senior Author}
\affil[1]{National University of Singapore}
\affil[2]{Partner Institution}

\paperlinks{%
  \ProjectLink{NUSPaper Project Website}{https://example.com/project}
  \DemoLink{Template Preview}{https://example.com/demo}
  \ArtifactLink{Supplementary Materials}{https://example.com/data}
  \SourceLink{GitHub Repository}{https://github.com/example/repo}
}

\correspondence{senior.author@example.edu}
\titlefootnote{* Equal contribution. \dag\ Corresponding author. This is an unofficial NUS-flavoured preprint style; use of the NUS logo should follow the university identity guidelines.}

\begin{document}
\maketitle

\begin{abstract}
This template provides a compact, lab-style LaTeX starting point for arXiv preprints, benchmark reports, and internal technical notes. It uses a clean first-page masthead, a centered high-contrast title block, compact author affiliations, a readable highlighted abstract, metadata links, an arXiv-style side stamp, and a bottom correspondence note. \textbf{The goal is to preserve the readability of a research paper while giving the document a recognizable house style before it is adapted to a venue-specific format.} The body pages remain conventional: single-column text, running headers, clean captions, theorem environments, code listings, and booktabs tables.
\end{abstract}

\keywords{LaTeX template; preprint; technical report; NUS; research communication}

\section{Introduction}

Research groups often need a paper format that is more polished than the default \texttt{article} class, but less restrictive than an official conference template. The common use case is a preprint, benchmark report, model card, dataset paper, or technical note that needs to look intentional on arXiv before it is adapted to a venue-specific format.

\begin{takeawaybox}
Use this template for public preprints and lab reports. For a formal conference submission, switch back to the venue-provided class and keep only the shared macros, figures, tables, and bibliography.
\end{takeawaybox}

The design philosophy is deliberately simple: a strong first page, readable body typography, restrained colour, and minimal custom syntax. The first page is meant to communicate the paper's identity, abstract, release links, and correspondence metadata at a glance, while subsequent pages behave like a conventional academic manuscript.

\paragraph{Design goals.}
The template optimizes for four properties:
\begin{enumerate}
  \item \textbf{Readable single-column layout}: better for arXiv, technical reports, and early drafts.
  \item \textbf{High-signal title page}: title, authors, affiliations, abstract, links, keywords, and footnotes are visible immediately.
  \item \textbf{Submission-safe structure}: content remains standard LaTeX and can be migrated to NeurIPS, ICLR, ICML, ACL, CVPR, or ACM classes.
  \item \textbf{Small style surface}: most visual changes live in \texttt{nuspaper.cls}, not in the manuscript body.
\end{enumerate}

\section{Template Structure}

The class exposes a small set of metadata commands in addition to the standard \LaTeX{} \verb|\title|, \verb|\author|, and \verb|\date| commands.

\begin{lstlisting}[language=TeX,caption={Minimal metadata block.}]
\documentclass{nuspaper}
\title{Your Paper Title}
\shorttitle{Running Header Title}
\date{2026-06-15}
\arxivstamp{arXiv:XXXX.XXXXX [cs.LG] 15 Jun 2026}
\author[1,*]{First Author}
\affil[1]{National University of Singapore}
\paperlinks{\FrontLink{Project Page}{Project Website}{https://example.com}}
\titlefootnote{* Equal contribution.}
\end{lstlisting}

\subsection{Recommended repository layout}

Keep the manuscript modular but not over-engineered. A practical structure is:

\begin{lstlisting}[language=bash,caption={Suggested file layout.}]
nuspaper-template/
|-- main.tex
|-- nuspaper.cls
|-- refs.bib
|-- sections/
|   |-- 01_introduction.tex
|   |-- 02_method.tex
|   `-- 03_experiments.tex
|-- figures/
|-- tables/
`-- latexmkrc
\end{lstlisting}

\subsection{Mathematical environments}

Standard theorem-like environments are included. For example:

\begin{definition}[House-style invariant]
A paper template is \emph{venue-compatible} if its scientific content can be moved to an official submission template without rewriting the mathematical notation, citations, labels, figures, or tables.
\end{definition}

\begin{proposition}
A small class file plus standard manuscript body files is easier to maintain than a large, paper-specific preamble copied across projects.
\end{proposition}

\begin{proof}
The class file centralizes typography and layout decisions. The manuscript body stays mostly semantic, so later migration to a venue template only requires replacing the class and adapting a small number of metadata commands.
\end{proof}

\section{Example Results}

Tables use \texttt{booktabs} by default. The class also defines light highlight helpers: \verb|\hlblue|, \verb|\hlorange|, and \verb|\hlred|.

\begin{table}[t]
\centering
\caption{Example benchmark table. Replace the synthetic numbers with real results. Highlighting is intentionally subtle so that the table remains printable.}
\label{tab:results}
\begin{tabular}{lrrrr}
\toprule
Model & Hard & Medium & Easy & Avg. score \\
\midrule
Baseline A      & \hlred{22.1} & 45.3 & 71.0 & 46.1 \\
Baseline B      & 31.5 & 52.4 & 78.8 & 54.2 \\
NUSPaper Demo   & \hlblue{44.2} & \hlblue{66.1} & \hlblue{84.0} & \hlorange{64.8} \\
\bottomrule
\end{tabular}
\end{table}

\Cref{tab:results} shows how a compact benchmark table looks in the template. Figure captions and table captions are ragged-right, small, and label-bold, which keeps long captions readable in single-column reports.

\subsection{Code snippets}

The bundled listing style is suitable for pseudocode, commands, and short implementation excerpts.

\begin{lstlisting}[language=Python,caption={Example analysis snippet.}]
def normalize_scores(scores):
    lo, hi = min(scores), max(scores)
    if hi == lo:
        return [0.0 for _ in scores]
    return [(x - lo) / (hi - lo) for x in scores]
\end{lstlisting}

\section{Conclusion}

This template is intended as a starting point rather than a final institutional style guide. The most useful next step is to replace the sample hero plot with the paper's central result, tune the colour accents to your lab identity, and add only the macros your group actually uses.


\bibliographystyle{plainnat}
\nocite{knuth1984texbook,example2026template}
\bibliography{refs}

\appendix
\section{Additional Notes}

\paragraph{Anonymous mode.}
Use \verb|\documentclass[anonymous]{nuspaper}| to suppress authors and affiliations on the title page.

\paragraph{Compact mode.}
Use \verb|\documentclass[compact]{nuspaper}| to reduce vertical spacing around section headings.

\paragraph{No stamp mode.}
Use \verb|\documentclass[nostamp]{nuspaper}| to hide the vertical arXiv-style stamp on the title page.

\paragraph{Paper size.}
The default is A4. Use \verb|\documentclass[letterpaper]{nuspaper}| for US Letter.


\end{document}
